\documentclass[10pt,a4paper]{amsart}
\usepackage[margin=19mm]{geometry}
\usepackage{amsmath,amssymb,mathtools}
\usepackage[scr=rsfs,cal=euler]{mathalfa}
\usepackage{xfrac}
\usepackage[unicode,hidelinks,bookmarks=false]{hyperref}
\newcommand{\ie}{i.e.\ }
\newcommand{\cf}{cf.\ }
\newcommand{\resp}{respectively\ }
\newcommand{\Subsection}{\subsection}
\providecommand{\nobreakdash}{\nobreak\hbox{-}}
\setlength{\emergencystretch}{2em}
\multlinegap0pt
\usepackage{bm}
\usepackage{bbm}
\usepackage{dsfont}
\usepackage{xspace}
\usepackage{cancel}
\usepackage{mathtools}
\usepackage{amsthm}
\makeatletter
\@ifundefined{theorem}{\newtheorem{theorem}{Theorem}}{}
\@ifundefined{proposition}{\newtheorem{proposition}[theorem]{Proposition}}{}
\makeatother
\newcommand{\dd}{\mathrm d}
\title[A generalization of the Fredenhagen-Haag derivation of Hawki]{A generalization of the Fredenhagen-Haag derivation of Hawking radiation for a class of Vaidya space-times}
\author{Felipe Dilho Alves}
\date{Source: arXiv:2608.03066v1 (CC BY 4.0)}
\begin{document}
\maketitle

\noindent\emph{Source and licence.} A generalization of the Fredenhagen-Haag derivation of Hawking radiation for a class of Vaidya space-times. Version arXiv:2608.03066v1 (CC BY 4.0), distributed under the Creative Commons Attribution 4.0 International licence, as recorded in the arXiv licence metadata for this exact version and shown on the abstract page https://arxiv.org/abs/2608.03066v1. Full rights notice: licenses/arxiv-2608.03066-CC-BY-4.0.txt.\par\smallskip

\noindent\textit{Editorial scope.} Counted unit: the Proposition whose source label is \texttt{eq:temporal-function} in the article (source0.tex lines 539--549, source label \texttt{eq:temporal-function}), delivered together with the article's own complete original proof of it (source0.tex lines 551--592, one \texttt{proof} environment of 42 lines). No mathematics is written, repaired or re-derived: statement and proof are the article's own text line for line. Required context retained uncounted: eq:unified-Vaidya (lines 488--495). Every definition, display and label of those lines is retained unchanged. Omitted from this package and not redistributed: 1--487, 496--538, 550--550, 593--3110. Nothing inside a retained excerpt is omitted or altered: every retained body is delivered exactly as the article's lines are written, so no omission receipt of any kind accompanies this package. Citations: the retained excerpts carry no citation command at all, so no bibliography entry of the article is transcribed and no reference is redistributed. Reference adaptation: none was needed and none was made; every reference inside the delivered unit resolves inside it, so no unresolved reference or citation is produced. Printed slips kept literal: no printed word, symbol, hypothesis or number of the counted statement or of its proof was altered. Layout and class: the article's own class is \texttt{article} with packages including fontenc, inputenc, lmodern, amsmath, amssymb, amsthm, mathtools, mathrsfs, microtype, geometry, booktabs, enumitem, xcolor, tikz; the wrapper is the lane's pinned \texttt{amsart} editorial wrapper at 10pt on A4, which restates the theorem-like environments the retained text needs and adds the article's own macro definitions that the retained text needs (\texttt{\textbackslash dd}), with extra wrapper packages (bm, bbm, dsfont, xspace, cancel, mathtools). The delivered numbering is the unit's own. Assets: no retained span contains a figure environment, a graphics-inclusion command, a PDF-page inclusion, a table or a TikZ picture; no image file is copied into this package, the package declares no asset dependency and no image file is redistributed with it. Licence: version arXiv:2608.03066v1 of this article is distributed under the Creative Commons Attribution 4.0 International licence, as recorded in the arXiv licence metadata for this exact version and shown on the abstract page https://arxiv.org/abs/2608.03066v1. A full rights notice is delivered at licenses/arxiv-2608.03066-CC-BY-4.0.txt. Characters: every retained excerpt is pure ASCII, so the delivered engine, XeLaTeX with its default Latin Modern text font, reports no missing character.
	\begin{equation}
		g_{\epsilon}
		=-C(w,r)\,\dd w^2+2\epsilon\,\dd w\,\dd r+r^2\dd\Omega^2,
		\qquad
		C(w,r)=1-\frac{2m(w)}{r},
		\qquad \epsilon\in\{+1,-1\}.
		\label{eq:unified-Vaidya}
	\end{equation}

	\begin{proposition}
		Let $m$ be smooth and nonnegative on the coordinate interval under
		consideration, and restrict to $r>0$.  For the metric
		\eqref{eq:unified-Vaidya}, the function
		\begin{equation}
			\tau_\epsilon=w-\epsilon r
			\label{eq:temporal-function}
		\end{equation}
		has an everywhere timelike gradient.  Consequently, every Vaidya coordinate
		region satisfying these assumptions is stably causal.
	\end{proposition}

	\begin{proof}
		The inverse of the $(w,r)$ part of the metric
		\eqref{eq:unified-Vaidya} is
		\begin{equation}
			g^{ww}=0,
			\qquad
			g^{wr}=g^{rw}=\epsilon,
			\qquad
			g^{rr}=C(w,r).
		\end{equation}
		Since
		\[
		\dd\tau_\epsilon=\dd w-\epsilon\,\dd r,
		\]
		we obtain
		\begin{align}
			g^{-1}(\dd\tau_\epsilon,\dd\tau_\epsilon)
			&=
			g^{ww}
			-2\epsilon g^{wr}
			+\epsilon^2g^{rr}\notag\\
			&=
			-2+C(w,r)\notag\\
			&=
			-1-\frac{2m(w)}r.
			\label{eq:temporal-calculation}
		\end{align}
		Because $m(w)\geq0$ and $r>0$,
		\[
		g^{-1}(\dd\tau_\epsilon,\dd\tau_\epsilon)
		\leq-1<0.
		\]
		Thus $\dd\tau_\epsilon$, or equivalently
		$\nabla\tau_\epsilon$, is everywhere timelike, so
		$\tau_\epsilon$ is a temporal function up to an overall choice of sign.
		
		The condition that a covector be timelike is open.  The light cones of $g$
		can therefore be widened slightly while keeping
		$\dd\tau_\epsilon$ timelike.  The function $\tau_\epsilon$ remains strictly
		monotone along every causal curve of the widened metric, and hence such a
		curve cannot be closed.  This proves stable causality.
	\end{proof}

\end{document}
